\section{Résultats}\label{sec:resultats}

\subsection{La baisse d'impôt ne fait entrer aucun capital neuf}

Le tableau~\ref{table:h1-capital} rapporte l'effet de la réforme de 2018 sur trois marges de
l'entrée de capital : la variation du capital souscrit en colonne~(1), la probabilité d'une hausse
du capital en colonne~(2) et le ratio de fonds propres en colonne~(3). Sur la première marge, le
coefficient de double différence vaut $-0{,}0034$, une grandeur économiquement nulle : rapportée au
niveau, elle signifie que le capital souscrit des sociétés françaises croît environ $0{,}34\%$ moins
vite que celui des pays comparables chaque année, un écart négligeable et de signe contraire à l'entrée de
capital qu'une baisse d'impôt attractive était censée provoquer.\footnote{$-0{,}0034$ est un
coefficient en quasi-logarithme, soit un écart de taux de croissance annuel d'environ
$-0{,}34\%$.} La permutation placebo place même la France au plus faible écart des six pays, avec
une valeur $p$ de $1{,}000$, c'est-à-dire que cinq pays non traités affichent un écart plus marqué
que le pays réformé. La probabilité d'une hausse du capital recule elle aussi très légèrement,
$-0{,}0049$, et le ratio de fonds propres, en colonne~(3), bouge de $+0{,}0474$ mais sans aucune
précision, la statistique de Student valant $+0{,}9$ et le placebo plaçant la France au deuxième rang
seulement. L'étude d'événement de la figure~\ref{fig:entree} confirme cette lecture : la trajectoire
de l'écart France contre les pays comparables est plate année après année avant 2018 et ne connaît aucune rupture
ensuite, les barres d'erreur couvrant zéro sur toute la fenêtre. L'hypothèse H1 est soutenue : la
baisse d'impôt ne fait pas entrer de capital neuf dans les fonds propres. Ce résultat n'est pas une
incertitude mais un nul estimé avec précision, cohérent avec la neutralité de
l'équation~\eqref{eq:neutralite} et avec l'absence d'effet réel documentée par \citet{yagan2015}
sur la baisse américaine de 2003.

\begin{figure}[ht]
  \centering
  \begin{tikzpicture}
    \begin{axis}[paperaxis,
      xlabel={Année}, ylabel={Écart France moins comparables (variation du log capital)},
      xtick={2014,2015,2016,2017,2018,2019,2020,2021,2022},
      xticklabel style={/pgf/number format/1000 sep=},
    ]
      \draw[zeroline, line width=0.6pt] (axis cs:2013.5,0) -- (axis cs:2022.5,0);
      % La France, en bleu, avec intervalle a 95 pour cent.
      \addplot[franceblue, mark=*, mark options={scale=0.8}, error bars/.cd,
        y dir=both, y explicit, error bar style={franceblue}]
        table[col sep=comma, x=year, y=coef, y error expr=1.96*\thisrow{std_error}]
        {figures/entree_capital.csv};
      % Reforme de 2018 : marqueur vert, la couleur de la baisse d'impot.
      \draw[baisse2018, dashed, line width=1pt]
        ({axis cs:2017.5,0}|-{rel axis cs:0,0}) -- ({axis cs:2017.5,0}|-{rel axis cs:0,1});
    \end{axis}
  \end{tikzpicture}
  \caption{Étude d'événement de l'entrée de capital autour de la réforme de 2018}
  \label{fig:entree}
  \fignote{Coefficients $\tau_j$ de l'équation~\eqref{eq:event} pour la variation du logarithme du
  capital souscrit, écart de la France aux cinq pays comparables l'année $j$, année de référence 2017 omise.
  Barres d'erreur à quatre-vingt-quinze pour cent sur l'erreur-type regroupée par pays, indicative.
  La ligne verte verticale marque la réforme de 2018. Un profil plat avant et après la réforme
  indique l'absence d'entrée de capital neuf.}
\end{figure}


{
\begin{table}[ht]
    \centering
    \caption{Entrée de capital après la réforme de 2018}\label{table:h1-capital}
    \tablenotes{Double différence de l'équation~\eqref{eq:did}, France contre les cinq pays comparables,
    effets fixes de société et d'année, contrôle de taille, période 2013 à 2022. La ligne
    \textit{France $\times$ après 2018} donne le coefficient $\beta$ ; la statistique de Student
    regroupée par pays est entre parenthèses. Les étoiles marquent la significativité d'après la
    valeur $p$ bilatérale associée à cette statistique : *** pour $p < 0{,}001$, ** pour
    $p < 0{,}01$, * pour $p < 0{,}05$. Le nombre de groupes étant de six, ces étoiles et la
    statistique de Student sont indicatives ; l'inférence du papier repose sur l'effet du contrôle synthétique et sur la
    valeur $p$ de permutation placebo, dont le plancher mécanique est $1/6 \approx 0{,}167$ ; une
    valeur de $1{,}000$ signifie que la France a l'écart le plus faible des six pays. Les variables
    continues sont winsorisées aux 1\textsuperscript{er} et 99\textsuperscript{e} centiles. Le ratio
    de fonds propres est calculé sur les années-sociétés à actif strictement positif. Les définitions
    sont dans la table~\ref{table:definitions-variables}.\\}
    \footnotesize
    \begin{tblr}{colspec={l*{3}{X[1,c]}}}
    \toprule
     & (1) & (2) & (3) \\
    Variables expliquées : & Variation du capital & Hausse du capital & Ratio de fonds propres \\
    \midrule
    France $\times$ après 2018 & $-0{,}0034^{***}$ & $-0{,}0049^{***}$ & $+0{,}0474$ \\
                               & $(-9{,}9)$  & $(-5{,}2)$  & $(+0{,}9)$  \\
     & & & \\
    Effet synthétique & $+0{,}0012$ & $-0{,}0059$ & $-0{,}0160$ \\
    Placebo $p$       & $1{,}000$   & $1{,}000$   & $0{,}833$   \\
     & & & \\
    Effets fixes société & Oui & Oui & Oui \\
    Effets fixes année   & Oui & Oui & Oui \\
    Contrôle de taille   & Oui & Oui & Oui \\
    Regroupement         & Pays & Pays & Pays \\
    Observations         & 16\,177\,643 & 16\,177\,643 & 16\,177\,643 \\
    \bottomrule
    \end{tblr}
\end{table}
}

\subsection{Absence de substitution par la dette et d'effet réel}

Le tableau~\ref{table:h2-substitution} rassemble les quatre marges de l'hypothèse de substitution :
le levier en colonne~(1), la croissance de la dette en colonne~(2), l'investissement en colonne~(3)
et la productivité du travail en colonne~(4). Aucune ne montre d'effet dans le sens attendu, et
aucune statistique de Student ne dépasse le seuil usuel. Le résultat le plus informatif tient à la
complémentarité comptable du financement : la part des fonds propres et la part des dettes dans le
bilan somment à l'unité, de sorte que le coefficient sur le levier, $-0{,}0472$, et celui sur le
ratio de fonds propres du tableau précédent, $+0{,}0474$, sont la même grandeur au signe près. S'il
y avait substitution, le point de fonds propres perdu se retrouverait en dette ; or les deux
coefficients sont nuls et imprécis, le placebo du levier plaçant la France au deuxième rang. Pour un
point de part de fonds propres en moins, on ne mesure donc pas un point de dette en plus, mais zéro.
L'investissement, en colonne~(3), affiche $+0{,}0683$ mais avec une statistique de $+1{,}5$ qui ne
franchit aucun seuil, et la productivité ne bouge pas. L'hypothèse H2 n'est pas soutenue au sens où
la substitution prédite par les anciennes vues n'apparaît pas ; conformément à
l'équation~\eqref{eq:neutralite}, les sociétés n'ont pas eu à remplacer des fonds distribués, faute
d'un surcroît d'investissement à financer. Cette absence d'effet réel prolonge le constat de
\citet{alstadsater2017} et de \citet{becker2013}, pour qui la fiscalité des versements déplace
l'allocation du capital sans en relever le niveau productif.

L'étude d'événement de la figure~\ref{fig:levier} confirme ce nul dans le temps. L'écart de levier
entre la France et ses pays comparables, un peu positif en début de période, reflue vers zéro avant
la réforme et ne remonte pas après 2018 ; il devient même légèrement négatif. Ce levier agrège
l'ensemble des dettes au bilan, et son absence de hausse signe l'absence de la substitution de la
dette aux fonds propres prédite par les anciennes vues. La dette financière garantie souscrite
pendant la crise sanitaire fait exception et monte en 2020 à 2022, mais elle est logée en trésorerie
plutôt que convertie en distribution ; ce mouvement est propre au Covid, non à la réforme, et nous
le caractérisons à la section~\ref{sec:robustesse}.

\begin{figure}[ht]
  \centering
  \begin{tikzpicture}
    \begin{axis}[paperaxis,
      xlabel={Année}, ylabel={Écart France moins comparables (levier)},
      xtick={2013,2014,2015,2016,2017,2018,2019,2020,2021,2022},
      xticklabel style={/pgf/number format/1000 sep=, rotate=45, anchor=east},
    ]
      \draw[zeroline, line width=0.6pt] (axis cs:2012.5,0) -- (axis cs:2022.5,0);
      % La France, en bleu, avec intervalle a 95 pour cent.
      \addplot[franceblue, mark=*, mark options={scale=0.8}, error bars/.cd,
        y dir=both, y explicit, error bar style={franceblue}]
        table[col sep=comma, x=year, y=coef, y error expr=1.96*\thisrow{std_error}]
        {figures/levier_event.csv};
      % Reforme de 2018 : marqueur vert, la couleur de la baisse d'impot.
      \draw[baisse2018, dashed, line width=1pt]
        ({axis cs:2017.5,0}|-{rel axis cs:0,0}) -- ({axis cs:2017.5,0}|-{rel axis cs:0,1});
    \end{axis}
  \end{tikzpicture}
  \caption{Étude d'événement du levier autour de la réforme de 2018}
  \label{fig:levier}
  \fignote{Coefficients $\tau_j$ de l'équation~\eqref{eq:event} pour le levier, total des dettes sur
  total de l'actif, écart de la France aux cinq pays comparables l'année $j$, année de référence 2017 omise.
  Barres d'erreur à quatre-vingt-quinze pour cent sur l'erreur-type regroupée par pays, indicative.
  La ligne verte verticale marque la réforme de 2018. Le levier français, un peu supérieur en début
  de période, reflue vers celui des pays comparables et ne remonte pas après la réforme. Ce levier
  agrège l'ensemble des dettes ; la hausse de la dette financière garantie pendant la crise
  sanitaire, logée en trésorerie, n'y apparaît pas et est caractérisée à la
  section~\ref{sec:robustesse}.}
\end{figure}


{
\begin{table}[ht]
    \centering
    \caption{Substitution par la dette et effets réels après 2018}\label{table:h2-substitution}
    \tablenotes{Double différence de l'équation~\eqref{eq:did}, mêmes spécification, période et
    conventions d'inférence que le tableau~\ref{table:h1-capital}. La statistique de Student
    regroupée par pays est entre parenthèses ; les étoiles marquent la significativité d'après la
    valeur $p$ bilatérale, *** pour $p < 0{,}001$, ** pour $p < 0{,}01$, * pour $p < 0{,}05$, et
    sont indicatives vu les six groupes. L'inférence repose sur l'effet synthétique et la valeur $p$
    de permutation placebo. Les définitions sont dans la table~\ref{table:definitions-variables}.\\}
    \footnotesize
    \begin{tblr}{colspec={l*{4}{X[1,c]}}}
    \toprule
     & (1) & (2) & (3) & (4) \\
    Variables expliquées : & Levier & Croissance dette & Investissement & Productivité \\
    \midrule
    France $\times$ après 2018 & $-0{,}0472$ & $+0{,}0302$ & $+0{,}0683$ & $+0{,}0049$ \\
                               & $(-0{,}9)$  & $(+0{,}2)$  & $(+1{,}5)$  & $(+0{,}4)$  \\
     & & & & \\
    Effet synthétique & $+0{,}0165$ & $+0{,}1763$ & $+0{,}0082$ & $-0{,}0896$ \\
    Placebo $p$       & $0{,}833$   & $0{,}333$   & $0{,}667$   & $0{,}500$   \\
     & & & & \\
    Effets fixes société & Oui & Oui & Oui & Oui \\
    Effets fixes année   & Oui & Oui & Oui & Oui \\
    Contrôle de taille   & Oui & Oui & Oui & Oui \\
    Regroupement         & Pays & Pays & Pays & Pays \\
    \bottomrule
    \end{tblr}
\end{table}
}

\subsection{La destination des fonds libérés : les ménages}

L'argent que la réforme ne dirige pas vers le capital des sociétés sort vers les ménages. Les
comptes nationaux montrent que les dividendes reçus par les ménages français augmentent de 32 \%
l'année même de la réforme, passant de 32\,592 à 43\,088 millions d'euros courants entre 2017 et
2018.\footnote{$43\,088 / 32\,592 - 1 = 0{,}322$ ; dividendes distribués reçus par les ménages
français, millions d'euros courants, 2017 et 2018, Eurostat.} La figure~\ref{fig:dividendes} place
la France parmi ses cinq pays comparables en indice base 100 en 2017 et isole nettement la rupture française :
l'indice français saute à 132 en 2018 quand la plus forte hausse parmi les pays comparables, celle des Pays-Bas,
atteint 114, de sorte que la France est le seul pays à rompre franchement l'année de sa réforme.
L'ampleur de ce saut, 32 points d'indice contre au plus 14 chez les pays comparables, mesure directement la
part des fonds qui sort du périmètre des sociétés. Du côté de l'emploi de ces fonds, les comptes
financiers des ménages montrent que le surcroît se loge en dépôts et en achats d'actions cotées sur
le marché secondaire, non en actions non cotées, c'est-à-dire non dans le capital des sociétés
privées, et l'indice des prix du logement français progresse moins que celui des pays comparables, ce qui ne
soutient pas un canal immobilier. L'argent reste donc dans l'épargne liquide et le portefeuille des
ménages ; il ne finance pas l'appareil productif, exactement l'écart entre réponse de la
distribution et réponse réelle que \citet{chetty2005} et \citet{yagan2015} relèvent sur la réforme
américaine de 2003.

Le contraste entre les deux versants du résultat se résume en une phrase. La
figure~\ref{fig:entree} ne montre aucune entrée de capital neuf dans les fonds propres des sociétés
françaises après la réforme de 2018, et le tableau~\ref{table:h2-substitution} aucun surcroît
d'investissement ni de dette ; la figure~\ref{fig:dividendes} montre au contraire une forte sortie
de dividendes vers les ménages, de 32 \% la seule année de la réforme. Autrement dit, après la
baisse d'impôt, rien n'entre dans le capital des entreprises et rien ne s'y investit, tandis qu'une
large part en sort vers les ménages. La baisse d'impôt n'a donc pas financé les entreprises ; elle a
alimenté la distribution.

\subsection{Une réponse asymétrique à la hausse et à la baisse d'impôt}

Nous plaçons enfin la hausse d'impôt de 2013 et la baisse de 2018 dans le même cadre afin de mesurer
l'asymétrie annoncée par l'équation~\eqref{eq:ajustement}. Sur la distribution, mesurée au niveau
agrégé des ménages, la semi-élasticité au choc de 2013 est fortement négative et celle au choc de
2018 positive mais plus petite. Sur une base lissée des trois années antérieures à chaque choc, qui
neutralise le report des versements d'une année isolée, la semi-élasticité vaut $-0{,}40$ pour la
hausse de 2013 et $+0{,}18$ pour la baisse de 2018 : la baisse ne défait donc que 45 \% de l'effet
de la hausse, un peu moins de la moitié.\footnote{$0{,}18 / 0{,}40 = 0{,}45$ ; semi-élasticités sur
base lissée des trois années antérieures, chocs de 2018 et de 2013.} Ce rapport ne tient pas au
choix de la fenêtre de référence : mesurée sur le flux brut d'une année sur l'autre, la
semi-élasticité vaut $-0{,}52$ en 2013 et $+0{,}23$ en 2018, soit le même rapport de 45
\%.\footnote{$0{,}23 / 0{,}52 = 0{,}45$ ; semi-élasticités sur le flux brut de l'année du choc, 2018
et 2013.}

Sur l'allocation du capital au niveau des sociétés privées, le tableau~\ref{table:h3-cliquet} montre
le même sens et chiffre l'asymétrie. La hausse de 2013 a réalloué du capital vers les bénéfices
retenus : la part des réserves dans l'actif monte de $+0{,}0942$ et la part des fonds propres de
$+0{,}1094$, l'une et l'autre nettement, avec des statistiques de Student de $+2{,}2$ et $+2{,}6$. La
baisse de 2018 ne ramène ces mêmes parts que de $+0{,}0332$ et $+0{,}0403$, des grandeurs trois fois
plus faibles et sans précision. La baisse ne défait donc qu'environ un tiers de la réallocation
opérée par la hausse, le capital retenu revenant dans un rapport de $1$ à $2{,}8$ sur les réserves et
de $1$ à $2{,}7$ sur les fonds propres.\footnote{$0{,}0942 / 0{,}0332 = 2{,}84$ pour la rétention,
$0{,}1094 / 0{,}0403 = 2{,}71$ pour les fonds propres ; coefficients de double différence aux chocs
de 2013 et de 2018.} La figure~\ref{fig:cliquet} résume cette asymétrie sur les trois marges, la
distribution, les réserves et les fonds propres, en portant côte à côte l'ampleur de chaque réponse :
la barre de la hausse de 2013 dépasse celle de la baisse de 2018 sur chacune. L'hypothèse H3 est
soutenue sur les deux marges, la distribution et le capital : conformément au cliquet d'ajustement
avec $\lambda^{+} > \lambda^{-}$, le capital part plus facilement qu'il ne revient.

\begin{figure}[ht]
  \centering
  \begin{tikzpicture}
    \begin{axis}[paperaxis,
      xlabel={Année}, ylabel={Dividendes reçus par les ménages (indice, 2017 $=100$)},
      xtick={2013,2014,2015,2016,2017,2018,2019,2020,2021,2022},
      xticklabel style={/pgf/number format/1000 sep=},
      legend pos=north west, legend cell align=left,
      ymin=40,
    ]
      % Cinq pays de comparaison, tous en gris : meme sens de couleur dans tout le papier.
      \addplot[controlgrey] table[col sep=comma, x=year, y=DE] {figures/dividendes_menages.csv};
      \addplot[controlgrey] table[col sep=comma, x=year, y=IT] {figures/dividendes_menages.csv};
      \addplot[controlgrey] table[col sep=comma, x=year, y=ES] {figures/dividendes_menages.csv};
      \addplot[controlgrey] table[col sep=comma, x=year, y=BE] {figures/dividendes_menages.csv};
      \addplot[controlgrey] table[col sep=comma, x=year, y=NL] {figures/dividendes_menages.csv};
      \addlegendentry{Pays de comparaison}
      % La France, toujours en bleu.
      \addplot[franceblue, line width=1.6pt, mark=*, mark options={scale=0.8}]
        table[col sep=comma, x=year, y=FR] {figures/dividendes_menages.csv};
      \addlegendentry{France}
      % Reforme de 2018 : marqueur vert, la couleur de la baisse d'impot.
      \draw[baisse2018, dashed, line width=1pt] (axis cs:2018,40) -- (axis cs:2018,200);
    \end{axis}
  \end{tikzpicture}
  \caption{Dividendes reçus par les ménages, France et cinq pays comparables, indice base 100 en 2017}
  \label{fig:dividendes}
  \fignote{Dividendes distribués reçus par les ménages, poste des comptes de secteur
  (ressources du secteur des ménages), millions d'euros courants, ramenés à un indice valant 100 en
  2017. En bleu la France, le pays traité ; en gris les cinq pays de comparaison, l'Allemagne,
  l'Italie, l'Espagne, la Belgique et les Pays-Bas. La ligne verte verticale marque la réforme de
  2018. Source : Eurostat, Office statistique de l'Union européenne.}
\end{figure}


\begin{figure}[ht]
  \centering
  \begin{tikzpicture}
    \begin{axis}[paperaxis,
      ybar, bar width=14pt,
      ylabel={Réponse à l'impôt (valeur absolue)},
      symbolic x coords={distribution,retention,fondspropres},
      xtick=data,
      xticklabels={Distribution, Réserves, Fonds propres},
      ymin=0,
      enlarge x limits=0.28,
      legend pos=north east, legend cell align=left,
      nodes near coords, nodes near coords style={font=\scriptsize},
      every node near coord/.append style={/pgf/number format/.cd, use comma, fixed, fixed zerofill, precision=2},
    ]
      % La hausse de 2013 en orange, la baisse de 2018 en vert : meme code que partout.
      \addplot[fill=hausse2013, draw=hausse2013!60!black]
        table[col sep=comma, x=marge, y=choc2013] {figures/cliquet.csv};
      \addlegendentry{Choc de hausse, 2013}
      \addplot[fill=baisse2018, draw=baisse2018!60!black]
        table[col sep=comma, x=marge, y=choc2018] {figures/cliquet.csv};
      \addlegendentry{Choc de baisse, 2018}
    \end{axis}
  \end{tikzpicture}
  \caption{Le cliquet fiscal : la réponse à la hausse de 2013 dépasse la réponse à la baisse de 2018}
  \label{fig:cliquet}
  \fignote{Ampleur de la réponse à chaque choc, en valeur absolue, sur trois marges. La distribution
  est la semi-élasticité des dividendes des ménages, sur base lissée des trois années antérieures au
  choc ; la rétention et les fonds propres sont les coefficients de double différence sur les parts
  de l'actif, au niveau des sociétés. Par convention de couleur, l'orange est la hausse d'impôt de
  2013 et le vert la baisse de 2018. Sur les trois marges, la réponse à la hausse de 2013 dépasse la
  réponse à la baisse de 2018 : le capital part plus facilement qu'il ne revient. La semi-élasticité
  de la distribution est négative en 2013 et positive en 2018 ; seules les amplitudes sont tracées.}
\end{figure}


{
\begin{table}[ht]
    \centering
    \caption{Le cliquet du capital : réallocation à la hausse de 2013 et à la baisse de 2018}
    \label{table:h3-cliquet}
    \tablenotes{Double différence de l'équation~\eqref{eq:did} estimée séparément autour de chaque
    choc, France contre les cinq pays comparables, effets fixes de société et d'année, contrôle de taille,
    panel étendu à 2010. La colonne \textit{Choc de 2013} prend l'année 2013 comme début de période
    traitée sur la fenêtre 2010 à 2016 ; la colonne \textit{Choc de 2018} prend 2018 sur la fenêtre
    2015 à 2022. La statistique de Student regroupée par pays est entre parenthèses ; les étoiles
    marquent la significativité d'après la valeur $p$ bilatérale, *** pour $p < 0{,}001$, ** pour
    $p < 0{,}01$, * pour $p < 0{,}05$, et sont indicatives vu les six groupes. La valeur $p$ de
    permutation placebo suit. Le dernier bloc donne la semi-élasticité de la distribution des dividendes des ménages, sur
    base lissée des trois années antérieures au choc, estimée au niveau agrégé. Les définitions sont
    dans la table~\ref{table:definitions-variables}.\\}
    \footnotesize
    \begin{tblr}{colspec={l*{2}{X[1,c]}}}
    \toprule
     & (1) & (2) \\
     & Choc de 2013 & Choc de 2018 \\
    \midrule
    \SetCell[c=3]{l}{\textit{Sociétés privées, part de l'actif}} & & \\
    Réserves sur actif & $+0{,}0942^{*}$ & $+0{,}0332$ \\
                       & $(+2{,}2)$  & $(+0{,}6)$  \\
    Fonds propres sur actif & $+0{,}1094^{**}$ & $+0{,}0403$ \\
                            & $(+2{,}6)$ & $(+0{,}8)$ \\
     & & \\
    Placebo $p$, réserves      & $0{,}667$ & $0{,}833$ \\
    Placebo $p$, fonds propres & $0{,}333$ & $0{,}833$ \\
     & & \\
    \SetCell[c=3]{l}{\textit{Ménages, distribution agrégée}} & & \\
    Semi-élasticité des dividendes & $-0{,}40$ & $+0{,}18$ \\
     & & \\
    Effets fixes société et année & Oui & Oui \\
    Contrôle de taille            & Oui & Oui \\
    Regroupement                  & Pays & Pays \\
    \bottomrule
    \end{tblr}
\end{table}
}
